%% DRAFT — Weight Optimization Evaluation + Performance Metrics (ADD to Results)
%% ACTION: Add as subsections in Section 6 (Results), after Baseline Pipeline Comparison.
%% Contains: held-out validation, LR baseline, tabular/graph ablation,
%%           graph extensions, synthesis, and computational performance.
%%
%% PLACEHOLDER CONVENTION: \textbf{[PENDING: description]} marks values to be filled
%% after the pipeline run completes.

\subsection{Held-Out Validation of Weight Optimization}
\label{sec:holdout_weight}

The scoring weights (Equations~1--2) were originally optimized on the full feature matrix and evaluated on the same data, producing an AUC of \textbf{[PENDING: full-data AUC from optimized\_weights.json]} recorded in \texttt{optimized\_weights.json}. Because Nelder-Mead explores the weight space against the very labels used for evaluation, that value is optimistically biased by construction. To determine whether the learned weights generalize, we re-ran the optimizer under a stratified calibration/evaluation split.

\subsubsection{Protocol}

A held-out evaluation wrapper loads the cached edge table via the same loader used by the feature audit, applies a stratified 50/50 split (seed~42), and partitions the \textbf{[PENDING: total edges]} masked edges into a calibration half (\textbf{[PENDING: cal edges]} edges, \textbf{[PENDING: cal red-team]} red-team) and an evaluation half (\textbf{[PENDING: eval edges]} edges, \textbf{[PENDING: eval red-team]} red-team). The optimizer is constructed on the calibration half only, runs Nelder-Mead to convergence on the same features used in the original pipeline (\textbf{[PENDING: feature list]}), and the resulting weights are applied to the evaluation half without further modification.

\subsubsection{Result: Optimization Generalizes}

Table~\ref{tab:holdout_optimizer} reports the calibration and evaluation AUCs.

\begin{table}[h]
\centering
\begin{tabular}{lc}
\hline
Quantity & Value \\
\hline
Calibration AUC & [PENDING] \\
Eval AUC (held-out) & \textbf{[PENDING]} \\
Calibration $-$ Eval gap & [PENDING] \\
Nelder-Mead iterations (calibration) & [PENDING] \\
\hline
\end{tabular}
\caption{Held-out validation of the Nelder-Mead weight optimizer. Calibration and evaluation AUCs agree to within \textbf{[PENDING: gap]}, indicating no measurable overfitting at $k$ free parameters over roughly \textbf{[PENDING: total edges]} edges.}
\label{tab:holdout_optimizer}
\end{table}

Calibration and evaluation AUCs agree to within \textbf{[PENDING: gap]}, well inside the variation expected from \textbf{[PENDING: red-team per split]} red-team edges per split. With $k$ free parameters against \textbf{[PENDING: total red-team]} red-team edges (approximately \textbf{[PENDING: params per positive]} parameters per positive sample), the hypothesis class is extremely under-parameterized, leaving essentially no capacity to overfit. The \textbf{[PENDING: full-data AUC]} figure in the original optimization trace is not an in-sample artifact; the same weights would land near \textbf{[PENDING: eval AUC]} on an unseen sample drawn from the same distribution. The AUC improvement from equal weights ($\approx$\textbf{[PENDING: equal-weight AUC]}) to optimized weights (\textbf{[PENDING: opt AUC]}) is therefore defensible with no held-out penalty.


\subsection{Logistic-Regression Baseline Comparison}
\label{sec:lr_baseline}

To test whether the bespoke Nelder-Mead optimizer adds value over a standard supervised linear method, we fit logistic regression (L2 penalty, balanced class weights) on the calibration half and evaluated on the held-out half, with identical feature preprocessing: the same percentile-rank transform applied to all non-binary features, followed by standardization on calibration-half statistics.

\begin{table}[h]
\centering
\begin{tabular}{lcccc}
\hline
Method & Cal AUC & Eval AUC & Gap & Wall time \\
\hline
Nelder-Mead optimizer   & [PENDING] & [PENDING] & [PENDING] & [PENDING] \\
Logistic regression (L2) & [PENDING] & \textbf{[PENDING]} & [PENDING] & [PENDING] \\
\hline
\end{tabular}
\caption{Nelder-Mead optimizer versus logistic regression on the same held-out split. The eval-AUC gap between methods is within single-seed noise.}
\label{tab:lr_comparison}
\end{table}

Logistic regression and the Nelder-Mead optimizer are essentially equivalent on this data (Table~\ref{tab:lr_comparison}). The specific choice of supervised linear method does not materially affect the result: Nelder-Mead, logistic regression, and most other supervised linear methods on these features all land near AUC~\textbf{[PENDING: AUC]}. The optimizer's value lies in producing interpretable weights aligned with the scoring formula, not in outperforming off-the-shelf alternatives.


\subsection{Tabular versus Graph-Derived Feature Ablation}
\label{sec:tabular_graph_ablation}

The Nelder-Mead-versus-LR comparison varies only the learner; both methods use the same features. To isolate whether the features themselves, and specifically the graph-derived ones, carry the discriminative signal, we ran a feature ablation that holds the learner constant (logistic regression, same held-out split, seed~42) and varies the feature set.

\subsubsection{Feature Grouping}

Features were split into two groups by whether they require graph topology to compute. \emph{Pure tabular} features are computable from per-edge attributes alone: authentication flags (\texttt{is\_ntlm}, \texttt{is\_network\_logon}, \texttt{is\_success\_auth}), edge rarity, protocol rarity, byte-per-packet ratio, duration z-score, and the unusual-destination-port flag --- 9 features total. \emph{Graph-derived} features require traversing the graph: source and destination node degrees (in, out, total), fan-out ratios, node-level temporal features (inter-arrival mean and std, burst score, active duration), source fan-out, destination fan-out ratio, and destination betweenness centrality --- 17 features total.

\subsubsection{Result: Graph-Derived Features Dominate}

\begin{table}[h]
\centering
\begin{tabular}{lccc}
\hline
Feature set & \# features & Cal AUC & Eval AUC \\
\hline
A. Pure tabular only    & 9  & [PENDING] & [PENDING] \\
B. Graph-derived only   & 17 & [PENDING] & \textbf{[PENDING]} \\
C. Combined (A $+$ B)  & 26 & [PENDING] & [PENDING] \\
\hline
\end{tabular}
\caption{Feature ablation under logistic regression. Adding graph-derived features on top of pure-tabular raises eval AUC by \textbf{[PENDING: $\Delta$ tab$\to$graph]}; adding pure-tabular on top of graph-derived raises eval AUC by only \textbf{[PENDING: $\Delta$ graph$\to$combined]}.}
\label{tab:tabular_graph}
\end{table}

Table~\ref{tab:tabular_graph} shows that graph-derived features carry roughly an order of magnitude more discriminative signal than pure-tabular features on this dataset. This result directly supports the graph-based methodology: the detection improvement comes from graph structure, not from the choice of supervised learner.


\subsection{Graph Feature Extensions}
\label{sec:graph_extensions}

The ablation in Section~\ref{sec:tabular_graph_ablation} shows that local graph features (degrees, fan-out, node-level temporal patterns) carry most of the signal. A natural follow-up is whether multi-hop graph reasoning, features that require traversing more than one hop, adds further improvement. We swept five candidate graph feature groups, added one at a time on top of the selected features under LR.

\begin{table}[h]
\centering
\begin{tabular}{lccc}
\hline
Feature group & Cal AUC & Eval AUC & $\Delta$ vs base \\
\hline
(base: selected features only)                         & [PENDING] & [PENDING] & --- \\
Standard PageRank                                & [PENDING] & [PENDING] & [PENDING] \\
Personalized PageRank (seed: \texttt{C17693})   & [PENDING] & [PENDING] & [PENDING] \\
k-core decomposition                            & [PENDING] & [PENDING] & [PENDING] \\
Louvain community                                & [PENDING] & [PENDING] & [PENDING] \\
Jaccard $+$ Adamic-Adar similarity              & [PENDING] & [PENDING] & [PENDING] \\
All five groups stacked                          & [PENDING] & [PENDING] & [PENDING] \\
\hline
\end{tabular}
\caption{Graph-feature sweep on top of the selected features under LR. Louvain community features include a cross-community flag and community-size features.}
\label{tab:graph_sweep}
\end{table}

\subsubsection{Personalized PageRank}

Personalized PageRank seeded at the known compromised host (\texttt{C17693}) adds far more than any other feature in the sweep (Table~\ref{tab:graph_sweep}). This is a real signal, but it carries a methodological caveat: the feature uses the identity of \texttt{C17693} as an input, and in a production deployment that identity would have to be discovered first. The result supports a \emph{conditional} claim: given a seed of confirmed-compromised hosts, structural propagation lifts eval AUC measurably. Carefully framed as a post-detection scoping mechanism, this is publishable; presented without qualification, it would overstate what the method can achieve under cold-start conditions.

\subsubsection{Label-Free Multi-Hop Features}

The label-free graph features, k-core decomposition and Louvain community membership, each add a modest but consistent improvement over the local-graph baseline. These features require no prior information about attackers and provide direct support for the claim that multi-hop graph structure carries detectable signal beyond local features. Standard PageRank and edge-similarity scores (Jaccard, Adamic-Adar) contribute essentially nothing on this dataset.


\subsection{Decomposing the Detection Improvement}
\label{sec:weight_synthesis}

The four results above isolate different axes of the detection pipeline. Table~\ref{tab:synthesis_axes} summarizes what each comparison tests.

\begin{table}[h]
\centering
\begin{tabular}{lp{2.2cm}p{2.0cm}p{2.2cm}}
\hline
Comparison & Held constant & What varies & Measures \\
\hline
Result~1 & Optimizer, features, split & Which half is evaluated & Overfitting? \emph{No.} \\
Result~2 & Features, split, supervision & Linear method & Optimizer choice? \emph{No measurable effect.} \\
Result~3 & Learner, supervision, split & Feature set & Graph features? \emph{Yes, \textbf{[PENDING: $\Delta$]}.} \\
Result~4 & Learner, base features, split & Graph feature group & Multi-hop features? \emph{Modestly, label-free.} \\
\hline
\end{tabular}
\caption{Summary of what each held-out comparison isolates.}
\label{tab:synthesis_axes}
\end{table}

Combining these results with the existing one-class baseline comparisons in Table~\ref{tab:methods_lanl} decomposes the improvement from na\"ive detection to the best achievable AUC into additive contributions (Table~\ref{tab:auc_budget}).

\begin{table}[h]
\centering
\begin{tabular}{lc}
\hline
Contribution & Approx.\ AUC $\Delta$ \\
\hline
Pure-tabular features under one-class learning & (lower bound) \\
Pure-tabular features under supervised LR      & [PENDING] \\
Graph-derived local features under supervised LR & [PENDING] \\
Multi-hop graph features (k-core, communities)  & [PENDING] \\
Known-attacker propagation (personalized PageRank) & [PENDING] (conditional) \\
\hline
\end{tabular}
\caption{Approximate AUC budget from the worst defensible baseline to the most aggressive supervised graph-aware system. The largest single addition is expected to be the graph-derived features themselves.}
\label{tab:auc_budget}
\end{table}

Three claims emerge from this evaluation with clear empirical backing. First, graph-derived features carry the bulk of the discriminative signal for lateral-movement detection on LANL-2015. Second, supervised learning is a separate, additive contribution from graph-derived features. Third, among supervised linear methods on the same features, the specific algorithm matters little: Nelder-Mead direct AUC maximization and logistic regression agree to within single-seed noise. The optimizer's value is in producing interpretable weights aligned with the scoring formula, not in outperforming standard alternatives.


\subsection{Computational Performance Metrics}
\label{sec:latency_throughput}

We report two complementary performance metrics per pipeline variant. \emph{Latency} is the total wall-clock time in seconds from the start of event streaming through graph construction, feature extraction, weight optimization, edge and path scoring, and threshold detection. Formally, $L = t_{\text{build}} + t_{\text{score}}$, where $t_{\text{build}}$ is the time to stream raw authentication and flow events into the graph and $t_{\text{score}}$ is the time for the full scoring pipeline (features, optimization, edges, paths, detection). \emph{Throughput} is the average number of events processed per second, computed as $T = N / L$, where $N$ is the total number of raw events (authentication plus flow records) streamed during graph construction. These metrics capture end-to-end pipeline cost rather than any single stage in isolation.
